%% kapa: exam=1 round=37 year=2026 period=1 form=A subject=economics subject_ko=경제학원론 seq=29 number=69
%% source: Q-Net 감정평가사 기출문제 (cst00309) artlSeq=5258924 file=raw/1008/5258924/2026년 제37회 감정평가사 1차 1교시.pdf
%% conversion: kordoc 4.24.51 -> tools/gen1.py (auto); review=proofread-sonnet; flags=none
\begin{question}{69}
총수요곡선 및 총공급곡선에 관한 설명으로 옳지 않은 것은? (단, 총수요곡선은 우하향, 총공급곡선은 우상향한다.)
\begin{choices}
  \item 수입 농산물 가격의 상승은 총공급곡선을 왼쪽으로 이동시킨다.
  \item 정부지출의 증가는 총수요곡선을 오른쪽으로 이동시킨다.
  \item 소비심리의 호전은 총수요곡선을 오른쪽으로 이동시킨다.
  \item 가계의 미래소득에 대한 낙관적인 전망은 총수요곡선을 오른쪽으로 이동시킨다.
  \item 통화량의 증가는 총수요곡선을 오른쪽으로, 총공급곡선을 왼쪽으로 이동시킨다.
\end{choices}
\end{question}
\answer{5}
